% sections_1_2.tex : title, abstract, Section 1 (Introduction) and Section 2 (Methods)
% for the LRC-BART revision (Phase 5b, 2026-09-09).
% Standalone body: \input this file into a wrapper carrying the preamble of
% paper/main.tex. Cites keys in paper/ref.bib and 05-writing/ref_additions.bib
% (merged in 05-writing/build/merged.bib).
%
% Theorem numbering. The main text states three results, Theorems 1 to 3, whose
% proofs are in the Web Appendix (03-theory/appendix_new.tex). The mapping is
%   Theorem 1  = Theorem A1 (bounded local influence) with Corollary A1;
%   Theorem 2  = Theorem A2(a),(b) and Theorem A3 (detection and the map), merged;
%   Theorem 3  = Theorem A4 (ESS calibration).
% With the paper's preamble (\newtheorem{theorem}{Theorem}[section]) these print as
% Theorems 2.1 to 2.3; dropping the [section] option in the preamble gives 1 to 3.
% Web-appendix cross references are written as "Web Appendix B" etc. in text.

\providecommand{\SIMRESULT}[1]{\textbf{[SIM:\ #1]}}

\title{Locally Robust Commensurate BART for Borrowing External Controls}
\author{}
\date{}
\maketitle

\begin{abstract}
Real-world data (RWD) can supplement the control arm of a randomized controlled trial (RCT), but the comparability of external and trial controls may vary with patient characteristics. Study-level meta-analytic predictive, commensurate, and power priors apply a common borrowing parameter across patients. We propose locally robust commensurate Bayesian additive regression trees (LRC-BART) to model the RCT control surface as a pooled BART surface plus a sparse discrepancy ensemble. A spike-and-slab commensurate prior on each discrepancy leaf allows borrowing to vary across covariate regions. The spike represents commensurability and the slab serves as the vague component of a robust meta-analytic predictive prior. Tree partitions determine the regions over which each component applies. All leaf marginals are exact, so the sampler targets the exact posterior. The prior effective sample size is defined at the estimand, the RCT-standardized control mean, and calibrated against the RWD posterior with an explicit ceiling. We prove that the RWD-induced shift of a leaf estimate is bounded and vanishes at a Gaussian rate under conflict, that the posterior probability of a small local discrepancy is consistent under fixed partitions and representability of the discrepancy, and that the calibration is monotone and well posed. A covariate-resolved map reports posterior compatibility at each patient profile; realized regional effective sizes summarize borrowing. Simulations with Gaussian and survival outcomes and a single-arm myeloma trial with a registry control illustrate the method.
\end{abstract}

\doublespacing

%==============================================================================
\section{Introduction}
\label{sec:intro}
%==============================================================================

Real-world data can supplement a randomized trial when collecting a sufficiently large control arm is difficult. Their usefulness depends on whether external outcomes represent the outcomes trial patients would have experienced under control. Trial selection, treatment practice and outcome ascertainment may vary with age, disease stage and prior therapy. External and trial controls may therefore be comparable for some patient profiles and differ substantially for others.

A common borrowing parameter cannot fully accommodate this variation. It may allow excessive borrowing where the sources disagree or discard useful information where they agree. We propose a Bayesian tree model that learns regions of covariate space from the data and allows borrowing to differ across them. We accompany the treatment-effect estimate with profile-specific summaries of compatibility and external information.

The established priors for historical borrowing act on a study-level parameter. The meta-analytic predictive (MAP) prior of \citet{Neuenschwander2010MAP} and \citet{schmidli2014robust} treats the control parameters of the historical studies and of the new study as exchangeable draws from a common population and takes the posterior predictive distribution of the new study's parameter as its prior; the robust MAP prior mixes that predictive component with a vague one so that a current study in conflict with the historical data can switch borrowing off. The commensurate prior of \citet{Hobbs2011,Hobbs2012} conditions the current parameter on the historical one, $\theta\given\theta_0\sim\Normal(\theta_0,1/\tau)$, and learns the commensurability precision $\tau$ from the data, with a spike-and-slab hyperprior so that the posterior can move between near-complete borrowing and near-none. The power prior of \citet{IbrahimChen2000} raises the historical likelihood to a power $a_0$, normalized by \citet{Duan2006} and \citet{Neuenschwander2009PP} so that $a_0$ can carry its own prior. \citet{RoverFriede2020,RoverFriede2025} show that with a single historical study, the case we treat here, the normal MAP, commensurate, and power-prior forms nearly coincide. In these study-level formulations, a common heterogeneity or discount scale determines borrowing without reference to individual covariates.

Locally robust commensurate BART (LRC-BART) models the two control sources jointly as a pooled Bayesian additive regression tree (BART) surface $f$ fitted to RCT and RWD controls alike, plus a sparse discrepancy ensemble $g$ that only the RCT controls inform, so that the RWD follow $f$ and the RCT controls follow $f+g$. Each leaf of $g$ carries a spike-and-slab prior whose spike is a commensurate component of small scale and whose slab is a vague component on the outcome scale, so the tree prior of $g$ decides, region by region, whether the RCT control surface is the RWD surface up to a calibrated discrepancy or a free local offset.

Our construction combines three established ideas. The decomposition into a pooled surface and a discrepancy ensemble multiplied by an indicator is the Bayesian causal forest of \citet{Hahn2020BCF}, with the source indicator in the place of the treatment indicator. The leaf prior is the commensurate prior of \citet{Hobbs2011} in spike-and-slab form, with the slab in the role of the vague component of \citet{schmidli2014robust}; conditioning on the pooled surface is the commensurate conditioning on the historical parameter, and since with one external source this is close to a single-study MAP prior \citep{RoverFriede2025}, we do not claim a methodological distinction on that basis. The calibration uses the expected local-information-ratio effective sample size of \citet{neuenschwander2020}. 

We combine these components by placing the spike-and-slab commensurate prior on the leaves of a discrepancy ensemble so that borrowing varies across covariate regions learned jointly with the surface, which a single ensemble with two means per leaf cannot do (Web Appendix~G.6). Every leaf marginal is in closed form, so the tree moves target the exact posterior and no surrogate likelihood is needed. The prior effective sample size is defined at the estimand, the RCT-standardized control mean, and calibrated against the RWD posterior; it has a ceiling equal to the information the RWD carries about that estimand, and a target can be stated as a fraction of the ceiling. We prove that the RWD-induced shift of a leaf estimate is bounded and vanishes at a Gaussian rate as the local conflict grows, that the posterior probability of a small local discrepancy is consistent under fixed partitions and representability of the discrepancy, and that the calibration is monotone and well posed. The posterior of $g$ then gives a covariate-resolved borrowing map, which reports at each patient profile the posterior probability that the local discrepancy is below a pre-specified threshold, together with a pointwise effective sample size. The model extends to right-censored outcomes through an accelerated failure time (AFT) likelihood with data augmentation, and to the single-arm setting in which no RCT control exists.

\subsection{Related work}
\label{sec:related}

A second group of methods lets the amount of borrowing vary with covariates, through propensity-score strata \citep{Wang2019PSPP,Chen2020PSCL}, a commensurate precision that depends on the covariates through kernel local regression \citep{Jin2023CAHB}, exchangeability over sources or over latent patient classes \citep{Kaizer2018MEM,Kotalik2021,Alt2024LEAP}, a global congruence statistic \citep{Jiang2023elastic}, or clustering of the patients on the covariates before the outcome model is fitted \citep{Bi2026SAMHC,Chandra2023}; Web Appendix~G.0 reviews them, and the closest parametric relative of our borrowing map is the covariate-adaptive hybrid borrowing of \citet{Jin2023CAHB}.

Among tree methods, \citet{ZhouJi2021} is the most directly related approach. They fit one BART ensemble \citep{Chipman2010BART} to RCT and external controls with the data source as a splitting covariate, so that trees pool the sources wherever they do not split on the source; borrowing is implicit and has no parameter, and their discussion proposes as future work a prior discouraging source splits together with a source offset whose scale controls borrowing. We model the proposed offset nonparametrically and give it a calibratable prior. \citet{Pan2022} apply BART to patient-level borrowing across subgroups. HE-BART \citep{Wundervald2022HEBART} places group-specific random effects inside each terminal node with a shared precision, the structural precedent for a hierarchical parameter inside a leaf, with the aim of mixed-effects prediction rather than controlling historical borrowing for a current study. flexBART \citep{Deshpande2025flexBART} lets a tree assign several levels of a categorical predictor to one branch, which with the source as predictor gives the implicit borrowing of \citet{ZhouJi2021} in a richer split space. None of these places a commensurate prior on a leaf, and among the covariate-adaptive methods only \citet{Jin2023CAHB} and \citet{Alt2024LEAP} report borrowing at the resolution of a patient profile, neither on a nonparametric outcome surface.

Section~\ref{sec:methods} sets out the model, the sampler, the effective sample size and its calibration, the borrowing map with its guarantees, and the single-arm variant. Section~\ref{sec:sim} reports simulation studies with Gaussian and survival outcomes, including scenarios in which the external data are compatible on part of the covariate space only, Section~\ref{sec:application} applies the method to a single-arm myeloma trial with a registry control, and Section~\ref{sec:discussion} concludes. Proofs are in the Web Appendix.

%==============================================================================
\section{Methods}
\label{sec:methods}
%==============================================================================

\subsection{Setting and estimands}
\label{sec:setting}

We observe an RCT with $N$ subjects, of whom $n_1$ are controls and $N-n_1$ are treated, and an external source of $n_2$ RWD controls. For subject $i$ we record the outcome $y_i$, the covariate vector $x_i\in\mathbb R^p$, and the source $g_i\in\{1,2\}$, with $g_i=1$ for an RCT control and $g_i=2$ for an RWD control; $D_i=\mathbf 1\{g_i=1\}$ indicates an RCT control. Continuous outcomes are modelled as Gaussian given the covariates. Survival outcomes are modelled on the log scale under a log-normal AFT model with right censoring, and $y_i$ then denotes the log survival time, observed or imputed as described in Section~\ref{sec:posterior}.

The estimands are standardized to the RCT covariate distribution. Let $f_T(x)$ be the mean outcome under treatment and $f_1(x)$ the mean outcome of an RCT control at covariate $x$. For Gaussian outcomes the average treatment effect is $\mathrm{ATE}=N^{-1}\sum_{i\le N}\{f_T(x_i)-f_1(x_i)\}$, the average over all $N$ RCT profiles, treated and control. For survival outcomes we report the ratio of population median survival times $q_{0.5}(\bar S_T)/q_{0.5}(\bar S_1)$, where $\bar S_a(t)=N^{-1}\sum_i S_a(t\given x_i)$ is the arm-$a$ survival function standardized to the RCT profiles, and the ratio of restricted mean survival times $\int_0^{t^\ast}\bar S_T(t)\,\dd t/\int_0^{t^\ast}\bar S_1(t)\,\dd t$ at a horizon $t^\ast$ inside the follow-up. The treatment surface $f_T$ is estimated from the RCT treated subjects alone with standard BART or AFT-BART. External observations inform the control surface $f_1$. We define effective sample size through the RCT-standardized control mean $\mu=N^{-1}\sum_{i\le N}f_1(x_i)$.

Both ensembles below are BART models \citep{Chipman2010BART}. A BART surface is a sum of $H$ regression trees, $\sum_{h=1}^H m(x;T_h,\theta_h)$, where $T_h$ is a binary tree whose internal nodes split on one covariate at one cutpoint and whose terminal nodes (leaves) $\ell=1,\dots,L_h$ carry means $\theta_{h\ell}$, and $m(x;T_h,\theta_h)$ returns the mean of the leaf to which $x$ is routed. The tree prior splits a node at depth $d$ with probability $\alpha(1+d)^{-\beta}$, chooses the splitting variable and cutpoint uniformly, and the leaf means are independent $\Normal(0,\lambda^2)$ with $\lambda=(\max y-\min y)/(2k\sqrt H)$ and $k=2$, the range rule, so that the sum of $H$ trees spans the outcome range with high prior probability.

\subsection{The model}
\label{sec:model}

The two control sources are modelled jointly as
\begin{equation}
y_i=f(x_i)+D_i\,g(x_i)+\varepsilon_i,\qquad \varepsilon_i\sim\Normal(0,\sigma^2_{g_i}),
\label{eq:model}
\end{equation}
Here $f(x)$ is the external-control response mean and $f_1(x)=f(x)+g(x)$ is the RCT-control response mean, with a separate residual variance for each source. Both surfaces are estimated jointly.

The pooled surface $f$ is a standard BART ensemble of $H_f=50$ trees with the default split prior $(\alpha_f,\beta_f)=(0.95,2)$ and leaf prior $\Normal(0,\lambda_f^2)$ with $\lambda_f$ from the range rule. The discrepancy surface $g$ is a small, shallow ensemble of $H_g=5$ trees with split prior $(\alpha_g,\beta_g)=(0.5,3)$, so that a $g$ tree is a stump a priori with probability one half and rarely deeper than two levels, and its leaves represent differences from the external-control response mean. The commensurate component centers $f_1(x)$ at $f(x)$. In particular, conditional on $f$, the discrepancy-tree partitions $T_g$ and the spike variance, if all $H_g$ leaves reached by $x$ are in the spike, then
\[
f_1(x)\given f,T_g,\{z_{h,\ell_h(x)}=1\}_{h=1}^{H_g},\tau_0^2
\sim\Normal\!\left(f(x),H_g\tau_0^2\right).
\]
Thus agreement means that the RCT-control response mean is close to the external-control response mean. We implement this prior with zero-centered leaves for $g(x)=f_1(x)-f(x)$:
\begin{equation}
\theta_{h\ell}\given z_{h\ell}\sim z_{h\ell}\,\Normal(0,\tau_0^2)+(1-z_{h\ell})\,\Normal(0,\tau_1^2),\qquad z_{h\ell}\iid\Ber(w),
\label{eq:leafprior}
\end{equation}
with $w\sim\mathrm{Beta}(1,1)$, $\tau_0^2\sim\text{scaled-Inv-}\chi^2(\nu_0,s_0^2)$ with $\nu_0=3$, and $\tau_1^2=(\max y-\min y)^2/(4k^2H_g)$ fixed by the range rule for an $H_g$-tree ensemble. We require $\tau_0^2<\tau_1^2$, Assumption A1 of the Web Appendix, and the sampler draws $\tau_0^2$ from its conditional truncated to $(0,\tau_1^2)$. The residual variances carry independent inverse-gamma priors. The spike, $z=1$, shrinks an individual discrepancy contribution toward zero and thereby shrinks the RCT response mean toward the external response mean. Section~\ref{sec:ess} calibrates its leaf scale $\tau_0$ through the effective sample size of the full ensemble. The slab, $z=0$, is the vague component of the robust MAP prior placed in a leaf of $g$: a free local discrepancy on the outcome scale, sized by the range rule. If the scale prior is left untruncated, integrating $z$ and $\tau_0^2$ gives the leaf marginal $w\,t_{\nu_0}(\theta;0,s_0^2)+(1-w)\Normal(\theta;0,\tau_1^2)$, The sampler uses the truncated scale prior, for which the spike marginal is not exactly Student-$t$. The local discrepancy $g(x)=f_1(x)-f(x)$ sums $H_g$ leaf contributions, one per tree. More generally, conditional on $f$, $T_g$, the reached-leaf indicators and both scales, $f_1(x)$ is normal with mean $f(x)$ and variance $\sum_{h=1}^{H_g}\{z_{h,\ell_h(x)}\tau_0^2+(1-z_{h,\ell_h(x)})\tau_1^2\}$. The slab increases the permitted departure from $f(x)$ while retaining the same prior center. These distributions describe the latent control mean; observation-level residual variance enters through \eqref{eq:model}.

Three special cases place the model among its comparators. With $w=1$ every leaf is in the spike and the model is a commensurate BART (C-BART) with a single commensurability scale $\tau_0^2$; with $w=1$ and $\tau_0^2\to0$ the discrepancy vanishes and the model is BART fitted to the pooled controls (BART-CP); with $w=0$ every leaf is in the slab and $g$ is a free nonparametric source offset regularized by the range rule, the tree analogue of BART with the source indicator as a covariate (BART-PP). LRC-BART with $w$ learned from the data moves between these per leaf.

The decomposition is what lets a regional shift be detected: in the alternative design, one ensemble with two means per leaf coupled through a common variance, the cheapest fit spreads a shift across all $H$ trees under the spike, and a prior-cost calculation (Web Appendix~G.6) shows that at $H=50$ no shift of clinical size flips an indicator, whereas in \eqref{eq:model} the shift lands in one or two leaves of a small $g$ ensemble after a single split, where the spike cannot hide it by spreading it. Identification is by support: where the RCT has no support $g$ is a prior draw and $f$ is the RWD surface, and where both sources have support $f$ and $g$ separate because $f$ is pinned by the RWD and $g$ is shrunk by the spike (Web Appendix~E).

\subsection{Posterior computation}
\label{sec:posterior}

We use Markov chain Monte Carlo with the backfitting scheme of \citet{Chipman2010BART}, one tree at a time with the others held fixed. The factors in each tree-move ratio have closed forms because conditionally on $(\tau_0^2,\tau_1^2,w)$ the leaf prior \eqref{eq:leafprior} is a finite Gaussian mixture. For a leaf of an $f$ tree the partial residuals $r_i=y_i-f_{-h}(x_i)-D_ig(x_i)$ come from both sources with their own variances; with $A=n_1/\sigma_1^2+n_2/\sigma_2^2$ and $B=S_1/\sigma_1^2+S_2/\sigma_2^2$, where $n_g$ and $S_g$ are the count and sum of the source-$g$ residuals in the leaf, the observation-level factors $\prod_i\phi(r_i;0,\sigma^2_{g_i})$ cancel from every tree-move ratio and the leaf factor is
\begin{equation}
M_f=(1+\lambda_f^2A)^{-1/2}\exp\Bigl\{\frac{\lambda_f^2B^2}{2(1+\lambda_f^2A)}\Bigr\},
\label{eq:Mf}
\end{equation}
the heteroscedastic form of the standard BART marginal. For a leaf of a $g$ tree the partial residuals $r_i=y_i-f(x_i)-g_{-h}(x_i)$ use only RCT controls in that leaf, with count $n_1$ and mean $\bar y_1$, and writing $v_1=\sigma_1^2/n_1$,
\begin{equation}
M_g=\frac{w\,\phi(\bar y_1;0,\tau_0^2+v_1)+(1-w)\,\phi(\bar y_1;0,\tau_1^2+v_1)}{\phi(\bar y_1;0,v_1)},
\label{eq:Mg}
\end{equation}
a two-component Gaussian mixture in the leaf mean, with $M_g=1$ in a leaf holding no RCT controls. Grow, prune, and change moves for both ensembles use the standard tree-prior and proposal ratios times the ratio of products of $M_f$ or $M_g$ over the affected leaves, a grow move contributing $M_LM_R/M_P$. Since these are the exact collapsed ratios, the tree kernels leave the exact posterior invariant with the leaf parameters and indicators integrated out, and the full sweep is a partially collapsed Gibbs sampler with the joint posterior as its invariant law (Lemma A1 and Proposition A1 of Web Appendix~A).

Conditional on the trees, we update the leaf parameters conjugately. An $f$ leaf draws $\theta\sim\Normal\bigl(\lambda_f^2B/(1+\lambda_f^2A),\ \lambda_f^2/(1+\lambda_f^2A)\bigr)$. A $g$ leaf is updated as a block: first
\begin{equation}
P(z=1\given\bar y_1)=\frac{w\,\phi(\bar y_1;0,\tau_0^2+v_1)}{w\,\phi(\bar y_1;0,\tau_0^2+v_1)+(1-w)\,\phi(\bar y_1;0,\tau_1^2+v_1)},
\label{eq:zpost}
\end{equation}
then $\theta\given z\sim\Normal\bigl(\tau_z^2\bar y_1/(\tau_z^2+v_1),\ \tau_z^2v_1/(\tau_z^2+v_1)\bigr)$ with $\tau_z^2$ equal to $\tau_0^2$ in the spike and $\tau_1^2$ in the slab. With $K_0$ spike leaves among the $L_g$ leaves of $g$ and $S_0$ the sum of their squared means, the hyperparameters update as $\tau_0^2\given\cdot\sim\text{scaled-Inv-}\chi^2\bigl(\nu_0+K_0,(\nu_0s_0^2+S_0)/(\nu_0+K_0)\bigr)$ truncated to $(0,\tau_1^2)$ and $w\given\cdot\sim\mathrm{Beta}(1+K_0,1+L_g-K_0)$, and the residual variances update from their inverse-gamma conditionals using the residuals $y_i-f(x_i)-D_ig(x_i)$ of each source. We fix $\tau_1^2$ by the range rule because slab leaves are few and the slab must allow discrepancies on the outcome scale. An $f$ leaf holding only one source is an ordinary BART leaf, and the minimum leaf size of five is enforced on each leaf's own data, pooled for $f$ and RCT-only for $g$. A sweep costs a standard BART on $n_1+n_2$ observations plus a small ensemble on $n_1$.

For right-censored survival outcomes we use the AFT model $\log y_i=f(x_i)+D_ig(x_i)+\varepsilon_i$ with data augmentation. At each iteration a censored subject's log survival time is imputed from $\Normal(f(x_i)+D_ig(x_i),\sigma_{g_i}^2)$ truncated to $(\log C_i,\infty)$, where $C_i$ is the censoring time, and treated as observed in the tree and parameter updates that follow; the treatment-arm AFT-BART is imputed the same way. The remaining tree and parameter updates are unchanged.

\subsection{Prior effective sample size and calibration}
\label{sec:ess}

The borrowing induced by a spike scale $s_0^2$ depends on the residual variance, leaf sizes and $H_g$. This makes the scale difficult to specify directly. We calibrate it through a prior effective sample size (ESS) defined at the estimand. Conditionally on the RWD, the prior for the RCT control surface is $f_1=f+g$ with $f\given\text{RWD}$ the posterior of a standard BART fitted to the RWD alone and $g$ drawn from its prior. For a scalar parameter with prior variance $V$ and unit Fisher information $1/\sigma_1^2$, the expected local-information-ratio ESS of \citet{neuenschwander2020} is $\sigma_1^2/V$, the number of RCT control observations carrying the same information as the prior, and we apply this to $\mu$. Write $V_\mu^f=\Var\bigl(N^{-1}\sum_{i\le N}f(x_i)\given\text{RWD}\bigr)$, which we estimate from posterior draws of the RWD-only BART evaluated at the RCT profiles, and $c_g=\E\bigl[N^{-2}\sum_{i,j}\mathbf 1\{x_i,x_j\text{ share a leaf}\}\bigr]$ under the $g$-tree prior, the leaf-sharing constant, which equals one for stumps and is computed once by Monte Carlo. Given the spike and slab states, $g$ adds $c_g\sum_h\tau^2_{z_h}$ to the prior variance of $\mu$. The ESS of the commensurate component, with every $g$ leaf in the spike, is
\begin{equation}
\mathrm{ESS}_0(s_0^2)=\sigma_1^2\,\E_{\tau_0^2\sim\text{scaled-Inv-}\chi^2(\nu_0,s_0^2)}\Bigl[\frac1{V_\mu^f+c_gH_g\tau_0^2}\Bigr],
\label{eq:ess0}
\end{equation}
a one-dimensional integral, and the prior-averaged ESS is
\[
\mathrm{ESS}(w,s_0^2)=\sum_{k=0}^{H_g}\binom{H_g}kw^k(1-w)^{H_g-k}\,\sigma_1^2\,\E\bigl[\{V_\mu^f+c_g(k\tau_0^2+(H_g-k)\tau_1^2)\}^{-1}\bigr].
\]
We calibrate $\mathrm{ESS}_0$ and report $w$ separately, as robust MAP reports the ESS of its informative component and the robust weight separately: conditional on $w$, the prior probability that every leaf covering a given profile is commensurate is $w^{H_g}$, so the averaged ESS is dominated by $w$ and cannot reach a target of the order of $n_2$ whatever $s_0^2$ is.

The prior ESS has a finite upper limit: as $s_0^2\to0$, $\mathrm{ESS}_0\to\sigma_1^2/V_\mu^f$, the information the RWD posterior carries about $\mu$ before any discrepancy is added. For a single tree with RCT leaf proportions equal to the RWD ones, $V_\mu^f=\sigma_2^2/n_2$ and the ceiling is $n_2\sigma_1^2/\sigma_2^2$ exactly, the RWD sample size in RCT units; under covariate shift it is lower, since RWD in regions without RCT patients carries no information about $\mu$ (Theorem~\ref{thm:ess}(d)). The regularization of $f$ itself contributes a term of order $\sigma_1^2/(H_f\lambda_f^2)$ that is negligible at $H_f=50$ and is not subtracted. A target $N_{\rm target}$ may be given as an absolute number of RCT controls, feasible only below the ceiling, or as a fraction of the ceiling, which is always feasible and which we recommend as the primary sensitivity ladder when the overlap is poor.

We first fit BART to the RWD and estimate $V_\mu^f$ from draws of the standardized surface. We divide the draws into $B$ blocks and rescale their variances $V^{(b)}$ to have mean $V_\mu^f$. For a grid $\mathcal S$ of spike scales, we use the Pr-rule
\begin{equation}
\hat s^2_{0\star}=\inf\bigl\{s_0^2\in\mathcal S:\widehat G(s_0^2)\ge q\bigr\},\qquad
\widehat G(s_0^2)=\frac1B\sum_{b=1}^B\mathbf 1\bigl\{\mathrm{ESS}_0(V^{(b)};s_0^2)\le N_{\rm target}\bigr\},
\label{eq:prrule}
\end{equation}
with $q=0.95$. This selects the most informative spike whose prior ESS is at most $N_{\rm target}$ in a fraction $q$ of the blocks. Targets at or above the whole-chain ceiling are flagged as capped. We report the prior $\mathrm{ESS}_0$ and the realized ESS, $\E_{\rm post}[\sigma_1^2/\{V_\mu^f+c_g\sum_h\tau^2_{z_h}\}]$. The pointwise quantity $\mathrm{ESS}(x)=\sigma_1^2/\{V_f(x)+H_g\tau_0^2\}$ describes local prior information, where $V_f(x)$ is the RWD posterior variance at $x$. Web Appendix~D gives the calibration mechanics.

\subsection{The borrowing map and its guarantees}
\label{sec:map}

We summarize compatibility over covariate space by the posterior probability that the local discrepancy is smaller than a threshold $c>0$ fixed in advance on the outcome scale,
\begin{equation}
\mathcal B_c(x)=P\bigl(|g(x)|<c\given\text{data}\bigr),
\label{eq:map}
\end{equation}
computed as the fraction of posterior draws, together with the posterior mean and interval of $g(x)$ and the pointwise ESS of Section~\ref{sec:ess}. We also retain the all-spike probability $B(x)$, the posterior probability that every leaf covering $x$ belongs to the spike, as a diagnostic; its prior scale is $w^{H_g}$ and, as Theorem~\ref{thm:map} explains, its limiting value need not identify the true compatibility state, whereas $\mathcal B_c$ is consistent under the stated conditions. In the simulation study $c=0.5$, one third of the RCT residual standard deviation for Gaussian outcomes and $0.5$ on the log-time scale for survival outcomes.

We give three theoretical results. The first two concern leaf-level inference and condition on the pooled surface $f$, on the partition of the $g$ trees, and on $(\sigma_1^2,\tau_0^2,\tau_1^2,w)$, and consider a $g$ leaf with $n$ RCT controls, residual mean $\bar y$, and $v=\sigma_1^2/n$. Write $\rho_j=\tau_j^2/(\tau_j^2+v)$ for $j\in\{0,1\}$, so that $\hat\theta_1=\rho_1\bar y$ is the slab estimate, close to the RCT-only residual mean, and $\hat\theta_0=\rho_0\bar y$ the spike-shrunk estimate, close to zero, that is, to the pooled surface. Let $p=P(z=0\given\bar y)$ be the posterior slab probability and
\begin{equation}
\kappa=\frac12\Bigl(\frac1{\tau_0^2+v}-\frac1{\tau_1^2+v}\Bigr),\qquad
a=\log\frac{1-w}{w}+\frac12\log\frac{\tau_0^2+v}{\tau_1^2+v},\qquad C=e^{-a}.
\label{eq:kappa}
\end{equation}

\begin{theorem}[Bounded local influence]\label{thm:influence}
Under Assumption A1 and the conditioning above, $\theta\given\bar y$ is the mixture $p\,\Normal(\hat\theta_1,\rho_1v)+(1-p)\,\Normal(\hat\theta_0,\rho_0v)$ with $\logit p=a+\kappa\bar y^2$, and the shift of the posterior mean relative to the slab estimate, $I(\bar y)=\E[\theta\given\bar y]-\hat\theta_1=-(1-p)(\rho_1-\rho_0)\bar y$, satisfies
\[
\sup_{\bar y}|I(\bar y)|\le\sqrt{2v\max\{\log C,\tfrac12\}},\qquad 1-p\le Ce^{-\kappa\bar y^2},
\]
so $I(\bar y)\to0$ at a Gaussian rate as $|\bar y|\to\infty$. With $w=1$ (C-BART) the shift is $-(\rho_1-\rho_0)\bar y$, linear and unbounded, and with $w=1$, $\tau_0^2\to0$ (BART-CP) it is $-\rho_1\bar y$.
\end{theorem}

Theorem~\ref{thm:influence} is Theorem A1 of Web Appendix~B. The commensurate component changes a leaf estimate relative to the RCT-only estimate by a bounded amount, of the order of one RCT standard error of the leaf, and this borrowing-induced shift vanishes as local conflict grows. The bound is on the borrowing-specific part; the slab estimate itself is shrunk toward the pooled surface by the factor $\rho_1$, the range-rule regularization any BART leaf applies, which is slight at ordinary leaf sizes. C-BART and BART-CP have no such bound, and their bias grows linearly with the conflict. Corollary A1 carries the bound to the estimand: at fixed partitions of the $g$ ensemble the RWD-induced shift of the standardized control mean relative to the free-offset model is bounded uniformly in the RCT data, whereas under C-BART and BART-CP it is a nonzero linear functional of the data and hence unbounded.

\begin{theorem}[Detection and the borrowing map]\label{thm:map}
Under Assumption A1 and the conditioning above, let the RCT residuals in the leaf be $d+\epsilon_i$ with $\epsilon_i\iid\Normal(0,\sigma_1^2)$, where $d$ is the true local discrepancy.
\begin{enumerate}[label=(\alph*),leftmargin=2.2em]
\item The spike probability $P(z=1\given\bar y)=\{1+\exp(a+\kappa\bar y^2)\}^{-1}$ is decreasing in $|\bar y|$, and when $\log\frac w{1-w}+\frac12\log\frac{\tau_1^2+v}{\tau_0^2+v}>0$ it equals one half at
\[
d_{1/2}=\Bigl[\frac{2(\tau_0^2+v)(\tau_1^2+v)}{\tau_1^2-\tau_0^2}\Bigl\{\log\frac w{1-w}+\frac12\log\frac{\tau_1^2+v}{\tau_0^2+v}\Bigr\}\Bigr]^{1/2},
\]
a few multiples of $(\tau_0^2+\sigma_1^2/n)^{1/2}$, so the probability that the leaf abstains at true shift $d$ is $\bar\Phi\{(d_{1/2}-d)/\sqrt v\}+\bar\Phi\{(d_{1/2}+d)/\sqrt v\}$, with $\bar\Phi$ the standard normal survival function.
\item For a single $g$ tree, $\mathcal B_c(x)=p\,\Psi_1(c)+(1-p)\,\Psi_0(c)$ with $\Psi_j(c)=P\{|\Normal(\rho_j\bar y,\rho_jv)|<c\}$, and as $n\to\infty$ with $|d|\ne c$, $\mathcal B_c(x)\to\mathbf 1\{|d|<c\}$ almost surely.
\item As $n\to\infty$ with $\tau_0^2$ fixed, the spike probability converges to a limit in $(0,1)$ that is at most $\frac{w\tau_1}{(1-w)\tau_0}\exp\{-\frac{d^2}2(\tau_0^{-2}-\tau_1^{-2})\}$ for $d\ne0$ and exceeds $w$ for $d=0$; it tends to zero only as $|d|/\tau_0\to\infty$, or for every fixed $d\ne0$ if $\tau_0^2\to0$ with $n$.
\end{enumerate}
\end{theorem}

Theorem~\ref{thm:map} combines Theorems A2 and A3 of Web Appendix~C. The ensemble extension requires fixed partitions and a discrepancy representable on their additive span. The leaf results also condition on $f$; they do not establish consistency under the full posterior over both ensembles. Further interpretation is given in Web Appendix~G.3.

\begin{theorem}[Monotone, well-posed calibration]\label{thm:ess}
Fix $\sigma_1^2>0$, $V>0$, and $c_gH_g>0$, and write $\mathrm{ESS}_0(V;s)$ for \eqref{eq:ess0} at $V_\mu^f=V$ and $s=s_0^2$.
\begin{enumerate}[label=(\alph*),leftmargin=2.2em]
\item $s\mapsto\mathrm{ESS}_0(V;s)$ is continuous, convex, and strictly decreasing on $(0,\infty)$, with limits $\sigma_1^2/V$ as $s\downarrow0$ and $0$ as $s\uparrow\infty$, under the untruncated calibration prior in \eqref{eq:ess0}. Under truncation to $(0,\tau_1^2)$, continuity and strict decrease remain, but the limit as $s\uparrow\infty$ is $\sigma_1^2/(V+c_gH_g\tau_1^2)>0$; convexity is asserted only for the untruncated law.
\item For each $V$ the sublevel set $\{s:\mathrm{ESS}_0(V;s)\le N_{\rm target}\}$ is a half-line $[s^\dagger(V),\infty)$, so the population rule $G(s)=P\{s^\dagger(V)\le s\}$ is a distribution function and $\{s:G(s)\ge q\}=[s_\star,\infty)$ has a unique attained infimum; $s_\star>0$ if and only if $P(N_{\rm target}<\sigma_1^2/V)>1-q$, and otherwise $s_\star=0$, the point-mass limit.
\item Under the stationary, aperiodic Harris-ergodic chain assumptions and fixed block-length regime specified in Web Appendix~D, $\widehat G(s)\to G(s)$ almost surely for each $s$, uniformly in $s$ when $G$ is continuous, and on a finite grid with $G\ne q$ at every grid point the empirical selection equals $s_\star$ for all $B$ large enough, almost surely; under the additional geometric-ergodicity and regularity conditions in that appendix, a central limit theorem includes the block-rescaling term. With a fixed number of blocks and increasing block length, the rule instead converges to the grid approximation of the whole-chain ESS root.
\item For a single $f$ tree with fixed partition, RCT profile proportions $a_\ell$ and RWD proportions $\pi_\ell$ over its leaves, $\lim_{s\downarrow0}\mathrm{ESS}_0(s)\ge n_2\sigma_1^2/\sigma_2^2\big/\sum_\ell a_\ell^2/\pi_\ell$, with equality in the flat-prior limit; the denominator is at least one, with equality if and only if $a=\pi$.
\end{enumerate}
\end{theorem}

Theorem~\ref{thm:ess} is Theorem A4 of Web Appendix~D. Parts (a) and (b) establish that the Pr-rule is well defined: the ESS decreases monotonically with the spike scale, so there is exactly one most informative admissible scale, and it is strictly positive exactly when the target is below the ceiling in enough Stage-1 blocks. A leaf-level ESS diverges as the scale shrinks and is feasible for any target; here the ceiling is finite and feasibility depends on the data and the target, which motivates specifying a fractional target. Part (c) gives conditions under which the Monte Carlo rule recovers the population selection. Part (d) relates the ceiling to the sample sizes: with matched covariate distributions the ceiling is the RWD sample size in RCT units, and covariate shift between the sources lowers it by the factor $\sum_\ell a_\ell^2/\pi_\ell$, which measures how unevenly the RCT profiles are spread over the leaves relative to the RWD.

\subsection{The single-arm variant}
\label{sec:singlearm}

When the trial has no control arm, $n_1=0$, every $g$ leaf holds no RCT controls, $M_g=1$, and the $f$ moves are exactly those of standard BART on the RWD. The control surface is then $f_1=f+g$ with $f$ the RWD surface and $g$ a prior draw, a zero-mean discrepancy process of pointwise variance $H_g\{w\tau_0^2+(1-w)\tau_1^2\}$ at a profile covered by $H_g$ distinct leaves, and $\tau_0^2$ and $w$ stay at their priors. The indicators receive no likelihood information, and slab draws only add prior noise of variance $(1-w)H_g\tau_1^2$ to the counterfactual control, so we report the $w=1$ model, the commensurate prior with the discrepancy carried by $g$ at the calibrated scale, and use $w<1$ as a sensitivity analysis. The calibration is unchanged: Stage 1 is the RWD-only fit, $\sigma_1^2$ is set to $\sigma_2^2$ in the absence of RCT outcomes, the ceiling is at most about $n_2$, and the target is stated as an absolute number of controls or a fraction of the ceiling. With $w=1$ the single-arm model is a hierarchically centered commensurate AFT-BART, $f_1(x)=f(x)+\sum_he_h$ with $e_h\sim\Normal(0,\tau_0^2)$, with the discrepancy defined on the partition of $g$ instead of that of $f$. Because the borrowing map is flat by construction here, the covariate-resolved display for a single-arm analysis is the local ESS map $\mathrm{ESS}(x)$, which describes variation in RWD support across trial profiles.
